\scanpage{27}

d'un $\unsure{\xi|\zeta|f}\in\mathrm{Ob}\,\mathcal F(X)$. À cause de l'unicité déjà acquise, la question est locale sur $X$, on peut supposer $X$ noethérien. Par suite, on peut supposer \deleted{que} (par récurrence noethérienne) que pour tt $Y'$ fermé $\subset Y$ et $\ne Y$, \deleted{\unsure{que l'on}} $\mathcal F$ satisfait $P_3$ rel. à $Y'$.

\deleted{Soit} Si \unsure{on} $Y=\emptyset$, il n'y a rien à prouver. Sinon, soit $x$ un pt maximal de $Y$.

Considérons $\xi'\mid X'_{(x)}$, par hypothèse il provient d'un $\xi_{0_{(x)}}\in\mathrm{Ob}\,\mathcal F_{(x)}(X_{(x)})$ lui-même provenant d'un $\xi_0\in\mathrm{Ob}\,\mathcal F(W)$, où $W$ est un voisinage \deleted{de} ouvert de $x$.

\deleted{On \unsure{a}} Considérons les restrictions de $\xi_0$ et $\xi'$ à $W\cap U=V$, par construction elles \deleted{coïncident} \inserted{deviennent isom} sur $V_{(x)}=\underset{\unsure{\nearrow|\to}}{X}'_{(x)}$, donc (par définition de $\mathcal F_{(x)}$) à condition de prendre $W$ assez petit, on peut supposer qu'ils sont isom. Choisissons un tel isom., on trouve donc un objet

$\xi'_1\in\mathrm{Ob}\,\mathcal F(W\cup U)$ prolongeant $\xi'$. Par l'hyp de récurrence, $\xi'_1$ provient d'un $\xi\in\mathcal F(X)$, on a fini.

% archive sheet number: 26
